% ==========================================================
% COURSE: DATA PERMISSIONS, SOURCES, TOOLS & CONNECTIONS
% ==========================================================
\documentclass[aspectratio=169, 11pt]{beamer}
\usetheme{Madrid}
\usecolortheme{whale}
\usefonttheme{professionalfonts}

\usepackage[utf8]{inputenc}
\usepackage{amsmath, amssymb}
\usepackage{booktabs}
\usepackage{xcolor}

\definecolor{primary}{RGB}{0, 82, 112}
\definecolor{accent}{RGB}{200, 83, 0}
\definecolor{softgray}{RGB}{245, 247, 250}

\setbeamercolor{title}{fg=white, bg=primary}
\setbeamercolor{frametitle}{fg=primary, bg=softgray}
\setbeamercolor{block title}{fg=white, bg=primary}
\setbeamercolor{block body}{fg=darkgray, bg=softgray}

\setbeamertemplate{navigation symbols}{}
\setbeamertemplate{footline}[frame number]

\title{\textbf{Data Permissions, Sources, Tools \& Connections}}
\subtitle{Enterprise Data Survival}
\author{Professional Development}
\date{}

\begin{document}
	
	% SLIDE 1: TITLE
	\begin{frame}
		\titlepage
		\centering
		\tiny{TOAD · Teradata · Vertica · Hive · ODBC · Oracle · PySpark · SAS}
	\end{frame}
	
	% SLIDE 2: COURSE OVERVIEW
	\begin{frame}{What this course covers}
		\tiny
		\begin{itemize}
			\item \textbf{Permissions} — prod vs dev vs uat, confidential levels
			\item \textbf{Sources} — where data lives, how to find it
			\item \textbf{Tools} — TOAD, Teradata, PySpark, PowerBroker
			\item \textbf{Connections} — ODBC, drivers, error handling
			\item \textbf{Queries} — SQL optimization, conversion to PySpark
			\item \textbf{SAS} — logistic regression, data cleaning
		\end{itemize}
	\end{frame}
	
	% SLIDE 3: DATA IS EVERYWHERE (AND MESSY)
	\begin{frame}{Enterprise data reality}
		\begin{block}{Different data sizes}
			\tiny
			\begin{itemize}
				\item \textbf{Data} — Excel, CSV, small tables
				\item \textbf{Big Data} — Terabytes, Hadoop, Hive
				\item \textbf{Very Big Data} — Petabytes, distributed clusters
			\end{itemize}
		\end{block}
		\vspace{0.03cm}
		\begin{block}{Different formats \& locations}
			\tiny
			\begin{itemize}
				\item Teradata, Vertica, Oracle, Hive, SQL Server
				\item On-premise vs cloud vs hybrid
				\item Prod / Dev / UAT / Sandbox environments
			\end{itemize}
		\end{block}
		\tiny\alert{You won't know where anything is. That's normal.}
	\end{frame}
	
	% SLIDE 4: PERMISSIONS — WHAT YOU CAN AND CANNOT DO
	\begin{frame}{Permissions hell}
		\begin{block}{Typical permission levels}
			\tiny
			\begin{itemize}
				\item \textbf{Prod} — live data, read-only usually
				\item \textbf{Dev} — test data, can write
				\item \textbf{UAT} — user acceptance testing
				\item \textbf{Sandbox} — play around
			\end{itemize}
		\end{block}
		\vspace{0.03cm}
		\begin{block}{The problem}
			\tiny
			\begin{itemize}
				\item You have Dev access but need Prod → error
				\item You have Prod access but not UAT → can't test
				\item You have no access at all → raise a ticket (wait 2 weeks)
			\end{itemize}
		\end{block}
		\tiny\alert{You will see "Table not found" when table exists. It's permissions.}
	\end{frame}
	
	% SLIDE 5: CONFIDENTIAL VS VERY CONFIDENTIAL
	\begin{frame}{Confidential data levels}
		\tiny
		\begin{columns}[T]
			\column{0.48\textwidth}
			\textbf{Confidential}
			\begin{itemize}
				\item Internal reports
				\item Non-public metrics
				\item Normal approval needed
			\end{itemize}
			\column{0.48\textwidth}
			\textbf{Very Confidential}
			\begin{itemize}
				\item PII (personally identifiable info)
				\item Customer data
				\item Legal/compliance review required
			\end{itemize}
		\end{columns}
		\vspace{0.1cm}
		\begin{block}{What this means for you}
			\tiny
			\begin{itemize}
				\item Different database schemas
				\item Different ODBC connections
				\item Extra approval forms
				\item "You are not authorized" errors
			\end{itemize}
		\end{block}
	\end{frame}
	
	% SLIDE 6: RAISING PERMISSIONS — THE PROCESS
	\begin{frame}{How to raise a permission request}
		\tiny
		\begin{enumerate}
			\item Get a new error (e.g., "Table not found")
			\item Search wiki for "permission" + database name
			\item Find the request form (SharePoint, ServiceNow, Jira)
			\item Fill it out (database, schema, table, read/write)
			\item Wait 2 days to 2 weeks
			\item Chase via email, Teams, Slack
			\item Eventually get access (or not)
		\end{enumerate}
		\vspace{0.05cm}
		\tiny\alert{No one will tell you this process. You find it yourself.}
	\end{frame}
	
	% SLIDE 7: PERMISSION ERRORS — WHAT THEY MEAN
	\begin{frame}{Common permission errors}
		\tiny
		\begin{tabular}{p{4cm} p{5cm}}
			\toprule
			\textbf{Error message} & \textbf{What it actually means} \\
			\midrule
			"Table not found" & Table exists, you can't see it \\
			"Database not found" & Wrong environment or no access \\
			"Access denied" & You need approval \\
			"ODBC connection failed" & Wrong driver or no network permission \\
			"User not authorized" & Ask your manager (who won't answer) \\
			\bottomrule
		\end{tabular}
	\end{frame}
	
	% SLIDE 8: CONNECTIONS OVERVIEW
	\begin{frame}{Connection types in enterprise}
		\tiny
		\begin{block}{ODBC (Open Database Connectivity)}
			\begin{itemize}
				\item \textbf{Manual ODBC} — you configure driver, DSN, server
				\item \textbf{Automatic ODBC} — pre-configured, just use it
			\end{itemize}
		\end{block}
		\begin{block}{Other connection methods}
			\tiny
			\begin{itemize}
				\item JDBC (Java)
				\item Native drivers (Teradata, Oracle)
				\item REST APIs
				\item SSH tunnels
			\end{itemize}
		\end{block}
	\end{frame}
	
	% SLIDE 9: TOAD CONNECTIONS
	\begin{frame}{TOAD connection screens}
		\tiny
		\begin{block}{TOAD Data Point (TDP)}
			\begin{itemize}
				\item Connect to Teradata, Oracle, SQL Server, Hadoop
				\item Need: Server name, port, database, username, password
				\item Authentication: LDAP, Windows, or database auth
			\end{itemize}
		\end{block}
		\vspace{0.03cm}
		\begin{block}{Common TOAD issues}
			\tiny
			\begin{itemize}
				\item "Unable to connect" → wrong driver or firewall
				\item "Invalid credentials" → password expired
				\item "Timeout" → network slow or query too big
			\end{itemize}
		\end{block}
	\end{frame}
	
	% SLIDE 10: TERADATA CONNECTIONS
	\begin{frame}{Teradata connections}
		\tiny
		\begin{block}{Via ODBC driver}
			\begin{itemize}
				\item Teradata ODBC driver installed (version matters)
				\item DSN configured with server IP and database
				\item Connection string: \texttt{DBCName=server; Database=schema;}
			\end{itemize}
		\end{block}
		\begin{block}{Common Teradata errors}
			\tiny
			\begin{itemize}
				\item "[Teradata] Driver] No server found" → wrong host
				\item "Logon failed" → wrong user/pass or expired
				\item "Query timeout" → query runs too long
			\end{itemize}
		\end{block}
	\end{frame}
	
	% SLIDE 11: VERTICA, HIVE, ORACLE
	\begin{frame}{Other databases — Vertica, Hive, Oracle}
		\tiny
		\begin{columns}[T]
			\column{0.3\textwidth}
			\textbf{Vertica}
			\begin{itemize}
				\item Columnar storage
				\item ODBC/JDBC
				\item High performance
			\end{itemize}
			\column{0.3\textwidth}
			\textbf{Hive}
			\begin{itemize}
				\item Hadoop layer
				\item SQL-like (HiveQL)
				\item Slow but big data
			\end{itemize}
			\column{0.3\textwidth}
			\textbf{Oracle}
			\begin{itemize}
				\item Traditional RDBMS
				\item TNS names file
				\item Enterprise standard
			\end{itemize}
		\end{columns}
		\vspace{0.1cm}
		\tiny\alert{Each has its own permission system. Each has its own error messages.}
	\end{frame}
	
	% SLIDE 12: CONNECTION ERRORS — COMPLETE LIST
	\begin{frame}{Connection errors you will see}
		\tiny
		\begin{itemize}
			\item "Unable to connect to server" — network down or firewall
			\item "Driver not found" — install ODBC driver
			\item "Authentication failed" — wrong password or expired
			\item "TNS could not resolve" — Oracle TNS misconfigured
			\item "SSL handshake failed" — certificate issue
			\item "Connection refused" — server not accepting
			\item "Timeout expired" — slow network or query
		\end{itemize}
		\vspace{0.05cm}
		\tiny\alert{You will see all of these. Probably in one day.}
	\end{frame}
	
	% SLIDE 13: TOOLS OVERVIEW
	\begin{frame}{Enterprise data tools}
		\tiny
		\begin{tabular}{p{3cm} p{6cm}}
			\toprule
			\textbf{Tool} & \textbf{Use} \\
			\midrule
			Teradata & Data warehouse, big SQL \\
			TOAD & Query editor, multiple DB connections \\
			PySpark & Big data processing, Python on clusters \\
			PowerBroker & Admin account escalation \\
			SAS & Statistics, logistic regression \\
			\bottomrule
		\end{tabular}
	\end{frame}
	
	% SLIDE 14: PYSPARK AND MASTER NODES
	\begin{frame}{PySpark on clusters}
		\tiny
		\begin{block}{Setup}
			\begin{itemize}
				\item Connect to master node (SSH)
				\item Submit job to cluster
				\item Data distributed across worker nodes
			\end{itemize}
		\end{block}
		\begin{block}{Common PySpark issues}
			\tiny
			\begin{itemize}
				\item "Master not found" — wrong IP
				\item "Out of memory" — executor too small
				\item "Stage failed" — data skew or missing partition
				\item "Permission denied" — wrong Linux user
			\end{itemize}
		\end{block}
	\end{frame}
	
	% SLIDE 15: QUERIES — LONG AND HARD
	\begin{frame}{Enterprise queries are not simple}
		\tiny
		\begin{block}{What you will see}
			\begin{itemize}
				\item 200+ line SQL queries
				\item 15+ joins across 10 tables
				\item Nested subqueries, CTEs, window functions
				\item No comments, no documentation
			\end{itemize}
		\end{block}
		\begin{block}{How to survive}
			\tiny
			\begin{itemize}
				\item Run small \texttt{SELECT COUNT(*)} first
				\item Add \texttt{LIMIT 10} or \texttt{SAMPLE}
				\item Comment as you go
				\item Break into smaller CTEs
			\end{itemize}
		\end{block}
	\end{frame}
	
	% SLIDE 16: JOINS — MULTIPLE TYPES
	\begin{frame}{Joins you will use}
		\tiny
		\begin{tabular}{p{3cm} p{5cm}}
			\toprule
			\textbf{Join type} & \textbf{When to use} \\
			\midrule
			INNER JOIN & Only matching rows \\
			LEFT JOIN & All rows from left table \\
			RIGHT JOIN & All rows from right table \\
			FULL OUTER JOIN & All rows from both \\
			CROSS JOIN & Cartesian product (dangerous) \\
			\bottomrule
		\end{tabular}
		\vspace{0.1cm}
		\tiny\alert{Mismatched keys, duplicate rows, and NULLs will break everything.}
	\end{frame}
	
	% SLIDE 17: AI FOR QUERIES — BE CAREFUL
	\begin{frame}{Using Copilot/AI for queries}
		\tiny
		\begin{block}{What AI can do}
			\begin{itemize}
				\item Convert SQL to PySpark
				\item Explain long queries
				\item Suggest optimizations
				\item Write simple queries from scratch
			\end{itemize}
		\end{block}
		\begin{block}{What AI will do wrong}
			\tiny
			\begin{itemize}
				\item Hallucinate columns that don't exist
				\item Ignore your database dialect
				\item Create inefficient joins
				\item Give confident wrong answers
			\end{itemize}
		\end{block}
		\tiny\alert{Never run AI-generated code without checking.}
	\end{frame}
	
	% SLIDE 18: QUERY OPTIMIZATION
	\begin{frame}{Making queries faster}
		\tiny
		\begin{enumerate}
			\item Use \texttt{WHERE} to filter early
			\item Avoid \texttt{SELECT *} — name columns
			\item Use \texttt{EXPLAIN} to see the plan
			\item Add indexes on join columns
			\item Partition large tables
			\item Avoid functions in \texttt{WHERE} clauses
			\item Use \texttt{UNION ALL} instead of \texttt{UNION} if possible
		\end{enumerate}
		\vspace{0.05cm}
		\tiny\alert{Slow queries will be killed by the DBA. You will get angry emails.}
	\end{frame}
	
	% SLIDE 19: SAS INTRODUCTION
	\begin{frame}{SAS in enterprise}
		\tiny
		\begin{block}{What SAS is used for}
			\begin{itemize}
				\item Logistic regression (credit risk, fraud)
				\item Data cleaning and transformation
				\item Statistical analysis
				\item Reporting (old school)
			\end{itemize}
		\end{block}
		\begin{block}{SAS code structure}
			\tiny
			\begin{itemize}
				\item \texttt{DATA} step — manipulate data
				\item \texttt{PROC} step — run procedures
				\item \texttt{PROC LOGISTIC} — regression
				\item \texttt{PROC FREQ} — frequency tables
			\end{itemize}
		\end{block}
	\end{frame}
	
	% SLIDE 20: READING SAS CODE
	\begin{frame}{How to read SAS code (when you don't know SAS)}
		\tiny
		\begin{block}{Example: Logistic regression}
			\scriptsize
			\texttt{proc logistic data=mydata;}\\
			\texttt{   class gender region;}\\
			\texttt{   model default(event='1') = age income debt / selection=stepwise;}\\
			\texttt{run;}
		\end{block}
		\vspace{0.05cm}
		\begin{block}{What it means}
			\tiny
			\begin{itemize}
				\item \texttt{proc logistic} — run logistic regression
				\item \texttt{class} — categorical variables
				\item \texttt{model} — dependent variable = independent vars
				\item \texttt{event='1'} — predict default = 1
			\end{itemize}
		\end{block}
	\end{frame}
	
	% SLIDE 21: SAS DATA CLEANING
	\begin{frame}{SAS for data cleaning}
		\tiny
		\begin{block}{Common SAS data steps}
			\scriptsize
			\texttt{data clean;}\\
			\texttt{   set raw;}\\
			\texttt{   if age = . then age = 0;}\\
			\texttt{   if income < 0 then delete;}\\
			\texttt{   format date yymmdd10.;}\\
			\texttt{run;}
		\end{block}
		\vspace{0.05cm}
		\begin{itemize}
			\item Handle missing values (.)
			\item Remove negative or invalid values
			\item Format dates properly
		\end{itemize}
	\end{frame}
	
	% SLIDE 22: IMPLEMENTATION TEAM AND FRAMEWORK
	\begin{frame}{How implementation team works}
		\tiny
		\begin{block}{What they do}
			\begin{itemize}
				\item Take model code from dev team
				\item Deploy to production environment
				\item Handle high resource requirements
				\item Monitor performance
			\end{itemize}
		\end{block}
		\begin{block}{Your interaction with them}
			\tiny
			\begin{itemize}
				\item Raise Jira tickets
				\item Wait (sometimes weeks)
				\item Get "deployment failed" errors
				\item Chase via Teams/Slack
			\end{itemize}
		\end{block}
	\end{frame}
	
	% SLIDE 23: HOW TO RUN A QUERY OR DATA REQUEST
	\begin{frame}{Step-by-step: running your first query}
		\tiny
		\begin{enumerate}
			\item Get connection details (wiki, email, teammate)
			\item Install ODBC driver (may need admin)
			\item Configure DSN (server, database)
			\item Open TOAD / Teradata / PySpark
			\item Write a simple \texttt{SELECT COUNT(*) FROM table}
			\item Run it
			\item Get an error
			\item Fix error (permissions, driver, typo)
			\item Repeat until it works
		\end{enumerate}
		\vspace{0.05cm}
		\tiny\alert{Step 7 is guaranteed. Step 8 is where you learn.}
	\end{frame}
	
	% SLIDE 24: SUMMARY — SURVIVING DATA HELL
	\begin{frame}{Enterprise data survival cheat sheet}
		\tiny
		\begin{tabular}{p{4cm} p{5cm}}
			\toprule
			\textbf{Problem} & \textbf{Solution} \\
			\midrule
			"Table not found" & Check permissions, then environment \\
			"Connection failed" & Check ODBC driver, firewall, VPN \\
			Query too slow & Add filters, limit, check indexes \\
			SAS won't run & Check library path, permissions \\
			No access to anything & Raise ticket, document, wait \\
			AI gave wrong code & Test on small sample first \\
			\bottomrule
		\end{tabular}
	\end{frame}
	
	% SLIDE 25: CLOSING
	\begin{frame}{You will survive data}
		\centering
		\tiny
		Every senior data person started exactly where you are now.
		
		No permissions. Broken connections. Unreadable queries. Silent teams.
		
		The difference?
		
		They kept trying. They documented what worked. They asked better questions.
		
		\vspace{0.2cm}
		\alert{Now you go do the same.}
		
		\vspace{0.2cm}
		\textit{Good luck. You will need it.}
	\end{frame}
	
\end{document}